\chapter{Sistema de predicción: agentes y meta-agente}
\label{chap:agentes}

El capítulo anterior estableció qué información existía en cada fecha. Éste explica cómo se
transforma esa información en una ordenación transversal de acciones. Es un capítulo de
arquitectura: describe el mecanismo, las alternativas que el catálogo ofrece en cada punto y las
restricciones que ninguna de ellas puede violar. No presenta todavía qué configuración se eligió ni
qué aprendió el sistema, y la omisión es deliberada.

La razón no es cosmética. Adelantar aquí los valores ganadores obligaría a justificar por qué ésos y
no otros, y esa justificación depende de una regla de decisión que todavía no se ha explicado; el
lector tendría que aceptarlos provisionalmente y volver atrás después. El documento sigue en su
lugar el orden lógico del argumento: primero el mecanismo y sus grados de libertad (este capítulo),
después la regla que elige entre ellos (Capítulo~\ref{chap:diseno-experimental}) y sólo entonces lo
que resultó de aplicarla (Capítulo~\ref{chap:resultados}).

Este capítulo describe la parte propiamente predictiva del sistema. La hipótesis no es que una sola
familia de ratios domine de forma estable todos los regímenes, sino que distintas vistas económicas
del universo aportan información con persistencia desigual. Por ello se entrenan cinco agentes
especializados ---\textit{quality}, \textit{value}, \textit{growth}, \textit{momentum} y
\textit{risk}--- y un combinador que aprende cómo ponderarlos sin conocer a priori cuál será el
más útil.

\section{Arquitectura y aprendizaje \textit{walk-forward}}

Cada agente recibe sólo los bloques de variables asignados a su especialidad. \textit{Quality}
resume rentabilidad, márgenes, eficiencia y solidez financiera; \textit{value} utiliza ratios de
precio y flujos; \textit{growth} mide aceleración y estabilidad de fundamentales; \textit{momentum}
recoge tendencias de precio y de fundamentales; y \textit{risk} representa riesgo de precio y
liquidez. Esta separación tiene dos fines: conserva interpretabilidad económica y evita que el
meta-agente confunda cinco copias del mismo predictor.

En cada fecha de snapshot se ajusta un LightGBM por agente sobre una ventana móvil de ocho años. El
catálogo declara la rejilla de hiperparámetros --- profundidad, número de estimadores, tasa de
aprendizaje y mínimo de observaciones por hoja --- y el protocolo del capítulo siguiente decide qué
valor toma cada uno. LightGBM permite representar interacciones no lineales sin imponer una forma
funcional única (Ke et al., 2017); el entrenamiento temporal evita que el modelo evalúe
observaciones que habría visto durante el ajuste.

Las cinco puntuaciones se convierten en rangos transversales y pasan al meta-agente, que estima una
Ridge positiva sobre cohortes trimestrales ya cerradas. Mientras no existan etiquetas cerradas
suficientes, el meta reparte por construcción pesos iguales de 0,20.

El catálogo ofrece dos variantes de esa combinación, y la elección entre ellas es una hipótesis de
diseño, no un detalle de implementación. \texttt{stacked\_rolling\_bounded} acota el peso de cada
agente entre 0,10 y 0,50: representa la hipótesis de que la arquitectura multi-agente sólo tiene
sentido si ningún especialista puede apagar a los demás, e impide por construcción que una ventaja
transitoria degenere el reparto. \texttt{stacked\_rolling\_free} no impone tope alguno: representa
la hipótesis contraria, que si un especialista domina de forma consistente lo correcto es dejar que
el procedimiento converja a él.

Las dos son defendibles a priori y miden cosas distintas. La acotada protege una propiedad
cualitativa del diseño --- que sigan siendo cinco agentes --- que ninguna métrica del protocolo
evalúa; la libre maximiza el criterio que el protocolo sí mide. Ésa es exactamente la clase de
disyuntiva que obliga a tener una regla de decisión declarada de antemano, y por eso la variante del
meta-agente figura en el catálogo cerrado como una variable más.

La Figura~\ref{fig:arquitectura} resume el recorrido completo, de las variables al ranking final.
Conviene fijarse en que el meta no ve las puntuaciones sino los \emph{rangos}: es lo que permite
combinar agentes cuyas escalas no son comparables entre sí.

\begin{figure}[H]
  \centering
  \begin{tikzpicture}[node distance=9mm and 12mm, >=Latex, every node/.style={font=\small}]
    \node[draw,rounded corners,fill=blue!8,minimum width=25mm,minimum height=10mm] (f) {Variables PIT};
    \node[draw,rounded corners,fill=teal!10,right=of f,minimum width=31mm,minimum height=22mm,align=center] (a) {Cinco agentes\\LightGBM};
    \node[draw,rounded corners,fill=orange!12,right=of a,minimum width=32mm,minimum height=16mm,align=center] (m) {Meta apilado\\sobre rangos};
    \node[draw,rounded corners,fill=green!10,right=of m,minimum width=26mm,minimum height=10mm,align=center] (s) {Ranking final};
    \draw[->,thick] (f) -- (a);
    \draw[->,thick] (a) -- (m);
    \draw[->,thick] (m) -- (s);
  \end{tikzpicture}
  \caption{Arquitectura predictiva. Las variables sólo están disponibles cuando ya se publicaron;
  los agentes se ajustan \textit{walk-forward} y el meta aprende con cohortes cuya etiqueta ya cerró.}
  \label{fig:arquitectura}
\end{figure}

\subsection{Qué predice el modelo}

El objetivo supervisado no es el retorno en euros ni una probabilidad de subida. Para cada fecha
$t$, se calcula el retorno total a doce meses $R_{i,t\rightarrow t+12}$ de cada acción $i$ cuya
etiqueta ya puede cerrarse y se transforma en percentil transversal:
\[
 y_{i,t}=\frac{\operatorname{rank}_{j\in\mathcal U_t}
 \left(R_{j,t\rightarrow t+12}\right)-1}{|\mathcal U_t|-1}.
\]
Así, $y=1$ identifica la mejor acción de su cohorte y $y=0$ la peor. Esta definición alinea
directamente el entrenamiento con el Rank-IC empleado para seleccionar: el sistema aprende
ordenación relativa, no trata de acertar el nivel absoluto de rentabilidad. También reduce la
dependencia de la escala de cada año, aunque no elimina la dependencia entre acciones que comparten
cohorte.

Cada agente $a$ estima una función no lineal
$\hat y^{(a)}_{i,t}=f_a(X^{(a)}_{i,t};\theta_{a,t})$, donde
$X^{(a)}_{i,t}$ contiene únicamente las variables permitidas a ese especialista y disponibles en
$t$. El subíndice temporal de $\theta$ importa: no existe un único árbol ajustado una vez sobre toda
la historia, sino una sucesión de modelos. Una predicción publicada en $t$ queda ligada a la versión
entrenada con la última información cerrada antes de $t$.

\subsection{Dos relojes que nunca se confunden}

La implementación mantiene separados el reloj de publicación de la señal y el reloj de cierre de la
etiqueta. En una fecha dada puede puntuarse una empresa con información observable, pero esa misma
fila no puede usarse para entrenar hasta que transcurran los doce meses necesarios para conocer su
resultado. La Figura~\ref{fig:dos-relojes} representa esa asimetría.

\begin{figure}[H]
  \centering
  \begin{tikzpicture}[>=Latex,every node/.style={font=\small},x=1.45cm,y=1cm]
    \draw[->,thick] (0,1.35) -- (8.4,1.35) node[right] {tiempo};
    \foreach \x/\lab in {0/$t-8a$,3.8/$t-12m$,5.5/$t$,8/$t+12m$}
      {\draw (\x,1.18)--(\x,1.52) node[above=2pt]{\lab};}
    \draw[very thick,blue!65] (0,0.85)--(3.8,0.85)
      node[midway,below]{labels cerrados para entrenar};
    \draw[very thick,teal!70!black] (0,0.25)--(5.5,0.25)
      node[midway,below]{variables observables para puntuar en $t$};
    \draw[very thick,orange!80!black,dashed] (5.5,-0.35)--(8,-0.35)
      node[midway,below]{retorno aún desconocido};
    \draw[->,thick] (5.5,0.15)--(5.5,-0.55);
  \end{tikzpicture}
  \caption{Separación entre el reloj de predicción y el de aprendizaje. La señal se emite en $t$;
  su etiqueta sólo entra en un ajuste posterior cuando el horizonte anual ha cerrado.}
  \label{fig:dos-relojes}
\end{figure}

Esta separación impide una fuga menos visible que usar un precio futuro: permitir que una fila
reciente, cuyas características sí se conocen pero cuyo retorno anual todavía no, intervenga en la
selección de variables o en el meta. El cierre de cohortes es por tanto una condición previa común
al entrenamiento, a la poda y a la calibración; no una corrección aplicada después.

\section{Qué aprende cada especialista}

La especialización no asigna una conclusión previa al modelo; asigna un vocabulario económico para
formular hipótesis. El agente \textit{quality} recibe ratios como ROE, ROIC, ROA, márgenes, rotación
de activos, liquidez y endeudamiento. Un ROIC alto puede indicar un uso eficiente del capital, pero
el modelo no recibe la regla ``ROIC alto implica comprar''; aprende una relación condicionada a la
historia disponible. \textit{Value} combina valoraciones relativas y flujos de caja; \textit{growth}
busca aceleración, persistencia y estabilidad de resultados; \textit{momentum} incorpora tendencia
de precio y de fundamentales; y \textit{risk} resume volatilidad, riesgo de precio y liquidez.

La amplitud de ese vocabulario es también una variable del catálogo. El preset \texttt{core} activa
un subconjunto reducido de bloques y el preset \texttt{all} los once disponibles; en ambos casos una
poda por estabilidad fuera de muestra limita cuántas variables conserva cada agente. La disyuntiva
es la habitual: un conjunto demasiado reducido amputa información relevante, y uno arbitrariamente
amplio aumenta la varianza y dificulta la interpretación. La poda forma parte del catálogo y opera
con datos ya observables; no es una selección manual posterior a ver el resultado.

\section{Entrenamiento y combinación de rangos}

En cada snapshot el modelo recibe sólo la historia permitida. Si la fecha es abril de 2020 y la
ventana es de ocho años, el conjunto comienza aproximadamente en abril de 2012; una fila de finales
de 2019 entra sólo si su etiqueta anual ya cerró. El ajuste se repite al avanzar el tiempo. Los
hiperparámetros que gobiernan la complejidad del árbol --- profundidad, número de estimadores, tasa
de aprendizaje y mínimo de observaciones por hoja --- controlan cuán específicas pueden ser las
interacciones aprendidas, y sus valores no se fijan aquí por criterio general sino con la regla del
Capítulo~\ref{chap:diseno-experimental}.

Los scores se convierten a rangos porcentuales dentro de cada cohorte. Esta operación evita que la
escala numérica de un agente domine al resto: si \textit{risk} produce valores entre 0,4 y 0,6 y
\textit{value} entre $-2$ y $3$, ambos pasan a expresar posición relativa en su sección transversal.
El meta combina después esos rangos, no sus unidades originales.

Antes del ajuste, los extremos se tratan con reglas declaradas y los valores ausentes se imputan
usando únicamente estadísticas obtenidas en el corte de entrenamiento. Una variable no se conserva
porque explique retrospectivamente toda la muestra: su disponibilidad y su contribución deben
sostenerse en los pliegues temporales permitidos. Con ello, la poda reduce dimensión sin convertir
el conjunto reservado en un selector informal. El diccionario completo de variables, direcciones,
bloques y agentes figura en el Apéndice~\ref{anexo:catalogo}.

Sean $z^{(a)}_{i,t}$ los cinco rangos producidos. Cuando existen suficientes cohortes cerradas, el
meta resuelve una regresión regularizada sobre predicciones verdaderamente fuera de muestra:
\[
 \widehat{\boldsymbol w}_t=
 \arg\min_{\boldsymbol w}
 \sum_{(i,s)\in\mathcal C_t}
 \left(y_{i,s}-w_0-\sum_{a=1}^{5}w_a z^{(a)}_{i,s}\right)^2
 +\lambda\lVert\boldsymbol w\rVert_2^2,
\]
donde $\mathcal C_t$ contiene sólo cohortes cerradas anteriores a $t$. Los coeficientes de agentes
se restringen a ser no negativos y después se normalizan para sumar uno. La variante acotada añade
$0{,}10\leq w_a\leq0{,}50$; la libre conserva únicamente no negatividad y normalización. El
intercepto sirve al ajuste, pero no altera la ordenación transversal final.

La regularización cumple aquí una función concreta: evita que agentes correlacionados intercambien
pesos extremos ante pequeñas perturbaciones. No garantiza, sin embargo, estabilidad económica. Si
un agente acumula más información ordinal que los demás, la versión libre puede concentrarse en él;
esa posibilidad es una hipótesis contrastada por el protocolo y se estudia después mediante pesos,
ablaciones y neutralización factorial.

Cuando $\mathcal C_t$ todavía no tiene la longitud mínima exigida, no se ajusta una regresión
frágil. El sistema utiliza el reparto igualitario y registra el estado de arranque. Esta conducta es
parte del contrato reproducible: dos ejecuciones con el mismo catálogo y la misma instantánea no
pueden decidir de manera distinta si ya existe historia suficiente.

\section{Ejemplo de combinación}

\begin{ejemplo}[De cinco rangos a un score]
Una empresa presenta rangos quality 0,80, value 0,55, growth 0,70, momentum 0,45 y risk 0,90. En el
arranque, los pesos iguales producen $(0,80+0,55+0,70+0,45+0,90)/5=0,68$: queda algo por encima de
la mediana del universo, y ninguna de sus cinco vistas pesa más que otra.

Si la evidencia ya cerrada indicase que \textit{risk} ordena mejor que los demás y el meta le
asignara 0,70, repartiendo 0,075 entre los cuatro restantes, el score pasaría a
$0{,}70\times0{,}90 + 0{,}075\times(0{,}80+0{,}55+0{,}70+0{,}45) = 0{,}82$.
\end{ejemplo}

Los rangos del ejemplo son inventados; lo que ilustra el cálculo es real. La empresa mejora de
posición, y conviene fijarse en por qué: no porque nadie haya decidido que el riesgo importa, sino
porque su fortaleza en la dimensión que mejor ha venido ordenando gana relevancia según datos ya
observados.

Aquí se ve con claridad qué está en juego entre las dos variantes del meta. Con la acotada, ese 0,70
sería imposible: el tope de 0,50 limitaría cuánto puede desplazarse el score hacia el especialista
dominante. Con la libre, nada impide que el peso siga creciendo mientras la evidencia lo respalde,
hasta el extremo de que el score final acabe siendo, en la práctica, el de un solo agente. Cuál de
las dos conductas es preferible no puede decidirse aquí, porque depende de qué muestren los datos, y
ése es justamente el trabajo del protocolo.

Conviene cerrar delimitando qué permite afirmar esta arquitectura. Los pesos del meta hacen
\emph{atribuible} el resultado a agentes concretos, pero no permiten inferir causalidad financiera:
que un especialista reciba peso alto no demuestra que su lectura económica sea la causa del retorno.
Puede reflejar interacciones entre variables, características del universo elegido o límites de los
datos disponibles. Por eso el trabajo no se detiene en los pesos y añade después estabilidad por
eras, neutralización por estilo y atribución local dentro de cada agente.

Con esto queda descrita la arquitectura completa: qué recibe cada agente, cómo se entrena, cómo se
combinan sus salidas y qué alternativas ofrece el catálogo en cada punto. Falta lo esencial para que
todo ello produzca un sistema concreto: la regla que elige entre esas alternativas sin mirar la
respuesta económica. El capítulo siguiente la describe.
